% !TEX root = ../main.tex
\corrigesChapitre{21}{Intensité sonore et atténuation}
\corrigesSouples

\begin{blocLivre}{cor-21-1}
\begin{corrige}{1}{Du niveau à l'intensité}
\demarche{$I=I_0\times10^{L/10}$ et $L=10\log(I/I_0)$.}
\begin{questions}
\item $I=10^{-12}\times10^{6}=\res{\qty{1,0e-6}{W.m^{-2}}}$.
\item $L=10\log\dfrac{\num{3.2e-5}}{\num{1e-12}}=\res{\qty{75}{dB}}$.
\item $I=10^{-12}\times10^{8{,}5}=\res{\qty{3,2e-4}{W.m^{-2}}}$.
\item $30$~dB d'écart correspondent à un rapport $10^{3}$ : l'intensité est \res{\text{1000 fois}} plus grande.
\end{questions}
\end{corrige}
\end{blocLivre}
\begin{blocLivre}{cor-21-2}
\begin{corrige}{2}{Un, deux, dix violons}
\demarche{Les intensités s'ajoutent : $N$ sources identiques donnent $L_1+10\log N$.}
\begin{questions}
\item $L=70+10\log2=\res{\qty{73}{dB}}$ : l'intensité a doublé, \res{\text{pas le niveau}}.
\item $L=70+10\log10=\res{\qty{80}{dB}}$.
\item $10\log N=15$, d'où $N=10^{1{,}5}=31{,}6$ : il faut \res{32} violons.
\item Les intensités audibles couvrent douze ordres de grandeur ; l'échelle logarithmique les ramène à des valeurs de
0 à \qty{120}{dB} et correspond mieux à la \res{\text{sensation auditive}}.
\end{questions}
\end{corrige}
\end{blocLivre}
\begin{blocLivre}{cor-21-3}
\begin{corrige}{3}{S'éloigner d'une enceinte}
\demarche{Source ponctuelle : $L_2=L_1-20\log(d_2/d_1)$ ; puissance $P=I\times4\pi d^2$.}
\begin{questions}
\item À \qty{4,0}{m} : $100-20\log2=\res{\qty{94}{dB}}$ ; à \qty{20}{m} : $100-20\log10=\res{\qty{80}{dB}}$.
\item Une atténuation \res{\text{géométrique}}.
\item $I=\qty{1,0e-2}{W.m^{-2}}$ à \qty{2,0}{m}, d'où $P=I\times4\pi d^2=\res{\qty{0,50}{W}}$.
\item $20\log\dfrac{d}{2{,}0}=30$, d'où $d=2{,}0\times10^{1{,}5}=\res{\qty{63}{m}}$.
\end{questions}
\end{corrige}
\end{blocLivre}
\begin{blocLivre}{cor-21-4}
\begin{corrige}{4}{Les protections auditives}
\demarche{Atténuation $A=L_1-L_2$ ; une atténuation de $A$ décibels divise l'intensité par $10^{A/10}$.}
\begin{questions}
\item Une atténuation \res{\text{par absorption}} : l'oreille reçoit $105-25=\res{\qty{80}{dB}}$.
\item $10^{-2{,}5}=\res{\num{3.2e-3}}$, soit 0,3~\% de l'intensité.
\item $94=85+3\times3$ : la durée est divisée par $2^3=8$, soit \res{\qty{1}{h}}.
\item À \qty{105}{dB}, la durée est divisée par $2^{20/3}\approx102$ : $\dfrac{8\times60}{102}\approx\res{\qty{4,7}{min}}$ seulement.
\end{questions}
\end{corrige}
\end{blocLivre}
\begin{blocLivre}{cor-21-5}
\begin{corrige}{5}{Logarithme et puissances de dix}
\demarche{$\log(10^n)=n$ ; $\log$ et $x\mapsto10^x$ sont réciproques.}
\begin{questions}
\item $\log1000=\res{3}$, $\log(10^{-12})=\res{-12}$ et $\log1=\res{0}$.
\item $\log2=\res{0{,}30}$ et $10^{0{,}5}=\res{3{,}16}$ : doubler l'intensité ajoute $10\times0{,}30=\qty{3}{dB}$.
\item $L/10=\log(I/I_0)$, donc $I/I_0=10^{L/10}$ : c'est la \res{\text{même relation}}.
\item De $10^{-12}$ à $10^{0}$ : \res{\text{12 ordres de grandeur}}, soit des niveaux de \qty{0}{dB} à \res{\qty{120}{dB}}.
\end{questions}
\end{corrige}
\end{blocLivre}
\begin{blocLivre}{cor-21-6}
\begin{corrige}{6}{Au pied de la scène}
\demarche{$L(d)=110-20\log(d/5{,}0)$.}
\begin{questions}
\item À \qty{20}{m} : $110-20\log4=\res{\qty{98}{dB}}$ ; à \qty{40}{m} : \res{\qty{91,9}{dB}}.
\item $20\log\dfrac{d}{5{,}0}=8$, d'où $d=\res{\qty{12,6}{m}}$.
\item $98-20=\res{\qty{78}{dB}}$.
\item L'air absorbe davantage les \res{\text{sons aigus}} : de loin, il reste surtout les graves.
\end{questions}
\end{corrige}
\end{blocLivre}
\begin{blocLivre}{cor-21-7}
\begin{corrige}{7}{Le brouhaha de la cantine}
\demarche{$N$ conversations identiques : $L=55+10\log N$.}
\begin{questions}
\item $L=55+10\log15=\res{\qty{66,8}{dB}}$.
\item $L=55+10\log30=\res{\qty{69,8}{dB}}$ : doubler le nombre de conversations n'ajoute que \res{\qty{3}{dB}}.
\item $10\log\dfrac{16}{15}=\res{\qty{0,28}{dB}}$ : une conversation de plus passe inaperçue.
\item Chacun élève la voix pour dominer le bruit ambiant, ce qui \res{\text{augmente ce bruit}} : tous parlent alors plus fort
(effet « cocktail »).
\end{questions}
\end{corrige}
\end{blocLivre}
\begin{blocLivre}{cor-21-8}
\begin{corrige}{8}{Programmer l'atténuation}
\demarche{Source ponctuelle : $I=\dfrac{P}{4\pi d^2}$ ; l'absorption retranche $\alpha\,d$ au niveau.}
\begin{questions}
\item \code{P} est la puissance sonore de la source (W), \code{I0} le seuil d'audibilité ($\unit{W.m^{-2}}$). La ligne traduit
la \res{\text{répartition de la puissance sur une sphère}} de surface $4\pi d^2$.
\item $99{,}0-93{,}0=93{,}0-87{,}0=\qty{6,0}{dB}$ ; puis $87{,}0-80{,}9=\qty{6,1}{dB}$ et $80{,}9-74{,}9=\qty{6,0}{dB}$, aux arrondis
près : \res{\qty{6}{dB}} par doublement.
\item \code{L = 10 * log10(I / I0) - 5 * d / 1000} : à \qty{16}{m}, le terme vaut \qty{0,08}{dB}, \res{\text{négligeable}}.
\item À \qty{1000}{m} : \res{\qty{39,0}{dB}} sans absorption, \res{\qty{34,0}{dB}} avec.
\end{questions}
\end{corrige}
\end{blocLivre}
\begin{blocLivre}{cor-21-9}
\begin{corrige}{9}{Choisir ses fenêtres}
\demarche{$L_\mathrm{int}=L_\mathrm{ext}-A$ ; $\qty{10}{dB}$ d'écart $=$ intensité divisée par 10.}
\begin{questions}
\item Une atténuation \res{\text{par absorption}} (et réflexion) : \res{\qty{47}{dB}} avec le simple vitrage, \res{\qty{37}{dB}} avec le
double vitrage.
\item $10^{(35-25)/10}=\res{10}$ fois plus faible.
\item $A\geqslant72-35=\res{\qty{37}{dB}}$.
\item $10^{-10/10}=\res{0{,}1}$ : 10~\% de l'intensité, cent fois plus qu'avec le double vitrage fermé.
\end{questions}
\end{corrige}
\end{blocLivre}
\begin{blocLivre}{cor-21-10}
\begin{corrige}{10}{Le grondement lointain du tonnerre}
\demarche{Atténuation géométrique $20\log(d_2/d_1)$, puis absorption $\alpha\times d$.}
\begin{questions}
\item $20\log\dfrac{1000}{1}=\res{\qty{60}{dB}}$.
\item À \qty{1}{km} : \res{\qty{37,2}{dB}} (\qty{500}{Hz}) ; \res{\qty{35,3}{dB}} (\qty{1000}{Hz}) ; \res{\qty{16}{dB}} (\qty{4000}{Hz}) ;
\res{\qty{-43}{dB}} (\qty{8000}{Hz}).
\item Les fréquences de 500 à \qty{4000}{Hz} restent audibles ; à \qty{8000}{Hz}, le niveau est \res{\text{sous le seuil}}
d'audibilité.
\item De près, toutes les fréquences parviennent, aigus compris : le tonnerre claque. De loin, l'air a absorbé les
\res{\text{aigus}} : il ne reste qu'un grondement grave.
\end{questions}
\end{corrige}
\end{blocLivre}
\begin{blocLivre}{cor-21-11}
\begin{corrige}{11}{Source ponctuelle ou source linéaire}
\demarche{$I=\dfrac{P}{2\pi d\,\ell}$, inversement proportionnelle à $d$ : $L_2=L_1-10\log(d_2/d_1)$.}
\begin{questions}
\item La puissance se répartit sur une surface $2\pi d\,\ell$ proportionnelle à $d$ : $I=\dfrac{P}{2\pi d\,\ell}$ est
\res{\text{inversement proportionnelle à }d}.
\item $I$ est divisée par 2 quand $d$ double : $\Delta L=-10\log2=\res{\qty{-3}{dB}}$.
\item $L=75-10\log4=\res{\qty{69}{dB}}$ ; pour une source ponctuelle, on aurait $75-20\log4=\qty{63}{dB}$.
\item $10\log\dfrac{d}{25}=15$, d'où $d=\res{\qty{791}{m}}$ : le niveau décroît lentement, d'où l'intérêt des murs
antibruit, qui ajoutent une forte atténuation près de la source.
\end{questions}
\end{corrige}
\end{blocLivre}
\begin{blocLivre}{cor-21-12}
\begin{corrige}{12}{Protéger le public d'un concert}
\demarche{Une enceinte : $110-20\log(d/4{,}0)\leqslant102$ ; deux enceintes : \qty{3}{dB} de plus, d'où $110+3-20\log(d/4{,}0)\leqslant102$.}
\begin{questions}
\item $20\log\dfrac{d}{4{,}0}\geqslant8$, soit $d\geqslant\res{\qty{10,0}{m}}$.
\item Les intensités s'ajoutent : $102+3=\res{\qty{105}{dB}}$, au-delà de la limite.
\item $20\log\dfrac{d}{4{,}0}\geqslant11$, soit $d\geqslant4{,}0\times10^{0{,}55}=\res{\qty{14,2}{m}}$.
\item Plusieurs petites enceintes réparties au plus près des spectateurs donnent un niveau \res{\text{plus homogène}} : on évite
des niveaux excessifs près de la scène pour obtenir un niveau suffisant au fond.
\end{questions}
\end{corrige}
\end{blocLivre}
\begin{blocLivre}{cor-21-13}
\begin{corrige}{13}{L'éolienne et ses voisins}
\demarche{$P=P_0\times10^{10{,}5}$, $I=\dfrac{P}{4\pi d^2}$, puis somme des intensités de l'éolienne et du bruit ambiant.}
\begin{questions}
\item $P=10^{-12}\times10^{10{,}5}=\res{\qty{0,032}{W}}$ ; à \qty{500}{m}, $I=\qty{1,0e-8}{W.m^{-2}}$ et $L=\res{\qty{40,0}{dB}}$.
\item $L_\mathrm{tot}=10\log(10^{3{,}5}+10^{4{,}0})=\res{\qty{41,2}{dB}}$ : émergence de \res{\qty{6,2}{dB}}, la réglementation
\res{\text{n'est pas respectée}}.
\item L'émergence vaut \qty{3}{dB} si $I_\text{éol}=I_0(10^{3{,}8}-10^{3{,}5})=\qty{3,15e-9}{W.m^{-2}}$, soit \qty{35,0}{dB} : il faut
$d=\sqrt{\dfrac{P}{4\pi I_\text{éol}}}=\res{\qty{894}{m}}$.
\item L'absorption réduit le niveau réel : le calcul est \res{\text{trop exigeant}}, la distance réellement nécessaire est plus
faible.
\end{questions}
\end{corrige}
\end{blocLivre}
\begin{blocLivre}{cor-21-14}
\begin{corrige}{14}{Choisir la fréquence d'une échographie}
\demarche{Atténuation $=0{,}5\times f\times2p$ ($f$ en MHz, profondeur $p$ en cm, aller et retour).}
\begin{questions}
\item $A=0{,}5\times5\times16=\res{\qty{40}{dB}}$.
\item $10^{-4}$ : seul un \res{\text{dix-millième}} de l'intensité revient.
\item $0{,}5\times7{,}5\times2p\leqslant60$ donne $p\leqslant\res{\qty{8,0}{cm}}$ ; à \qty{3,5}{MHz}, $p\leqslant\res{\qty{17,1}{cm}}$.
\item Thyroïde, peu profonde : une fréquence \res{\text{élevée}} (\qty{7,5}{MHz} ou plus), pour une image fine. Foie : une
fréquence \res{\text{basse}} (\qty{3,5}{MHz}), seule capable d'atteindre cette profondeur.
\end{questions}
\end{corrige}
\end{blocLivre}
\begin{blocLivre}{cor-21-15}
\begin{corrige}{15}{La dose de bruit quotidienne}
\demarche{L'énergie sonore reçue est proportionnelle à $I\times t$, avec $I\propto10^{L/10}$.}
\begin{questions}
\item L'énergie reçue est proportionnelle à l'\res{\text{intensité}} et à la durée : on ajoute les produits $I_i\,t_i$, et non
les niveaux.
\item $L_\mathrm{moy}=10\log\dfrac{10^{9{,}4}+3\times10^{8{,}8}}{8}=\res{\qty{87,4}{dB}}$ : au-delà de \qty{85}{dB}, les protections sont
\res{\text{obligatoires}}.
\item Les niveaux deviennent \qty{79}{dB} et \qty{73}{dB} : $L_\mathrm{moy}=\res{\qty{72,4}{dB}}$.
\item $1\times10^{9{,}4}=\num{2.51e9}$ et $8\times10^{8{,}5}=\num{2.53e9}$ : les deux expositions transportent la
\res{\text{même énergie}}, aux arrondis près.
\end{questions}
\end{corrige}
\end{blocLivre}
